The rapid ascendancy of social media platforms has radically restructured the architecture of modern human communication, shifting public discourse from local physical spaces to global digital environments. While these platforms were initially championed as instruments of democratic liberation and social connection, they are increasingly scrutinized for their role in fostering severe ideological and affective polarization \cite{Bail2018, Levy2021}. In these virtual arenas, disagreement is frequently amplified, leading to a fragmented public square characterized by echo chambers, mutual distrust, and hostile interactions. 

For social media firms, systemic polarization is not merely a regrettable social externality; it represents a profound threat to their core economic models. When digital environments become excessively hostile, users experience cumulative cognitive fatigue and emotional distress, leading to permanent exit from the platform---a phenomenon described in classic economic theory by \cite{Hirschman1970} as the ``exit'' option. Under severe polarization, user churn accelerates, directly reducing the active user base and eroding the platform's primary source of monetization: ad impressions. Furthermore, heightened toxicity and extreme discourse trigger a ``brand safety penalty'' \cite{Guess2023}. Advertisers, highly sensitive to reputational risks, increasingly restrict or withdraw their ad spend from platforms that fail to maintain a civil and safe environment. Consequently, social media corporations face a critical imperative to deploy algorithmic moderation strategies to maintain user retention and stabilize advertising revenue.

To combat this toxic churn, platforms have turned to algorithmic interventions. Historically, moderation focused on reactive ex-post content removal or account suspensions. More recently, however, platforms have begun experimenting with structural feed interventions, such as ``bridging algorithms'' \cite{Guess2023}. These recommendation systems seek to identify and suppress confrontational, cross-cutting content that is highly likely to trigger negative cognitive responses and aggressive exchanges. From a psychological standpoint, such interventions target the ``latitude of rejection'' \cite{Sherif1961}---the zone of opinions so divergent from an individual's prior beliefs that exposure causes cognitive backlash and opinion repulsion (the backfire effect) rather than assimilation. By suppressively filtering these hostile, backfire-inducing interactions, bridging algorithms seek to lower user frustration, minimize churn, and maintain a brand-safe digital ecosystem.

Yet, this study uncovers a profound and counter-intuitive tension between platform economics and systemic social dynamics, which we define as the \textit{Polarization Paradox}. While bridging algorithms are highly effective at achieving their immediate goals---successfully curbing user frustration, reducing churn, and preserving platform revenue---they simultaneously entrench and exacerbate systemic polarization. We argue that by shielding individuals from disagreeable, cross-cutting interactions, bridging algorithms prevent the constructive tension necessary for the cultivation of interpersonal tolerance and mutual trust.

We formalize this thesis by introducing a novel agent-based simulation framework: the Dynamic-BABE (Biased Assimilation \& Behavioral Entrenchment) model. Our framework extends classical models of opinion dynamics, such as the DeGroot consensus model \cite{DeGroot1974} and the Hegselmann-Krause bounded confidence model \cite{Hegselmann2002}, by integrating the entrenchment-based biased assimilation formulation of Chen et al. \cite{Chen2021} with Social Judgment Theory \cite{Sherif1961, Jager2005} and a co-evolving network of dyadic interpersonal trust. Crucially, trust in our model is not a static network parameter but a dynamic, dyadic property that co-evolves with opinion interactions. Consistent with empirical findings from risk and trust research \cite{Slovic1993}, trust is fragile: it builds slowly upon successful consensus-seeking interactions but erodes rapidly when confrontational, backfire-inducing exchanges occur.

Through a rigorous 2$\times$2 factorial experiment crossing the activation of the bridging algorithm and dynamic trust, we demonstrate that when dyadic trust is fragile, the suppression of confrontational cross-cutting interactions deprives agents of the opportunity to build mutual tolerance. When the bridging algorithm is active, agents are effectively insulated from their ideological opponents. While this dampens immediate conflict and eliminates the frustration that triggers churn, it also prevents the successful, low-friction interactions that allow trust to rebuild across ideological lines. Over time, relations that are not actively maintained wither due to passive relational decay. The network consequently undergoes \textit{trust segregation}---a state where trust is highly concentrated within homogeneous opinion clusters, while trust between opposing groups completely collapses. Ultimately, this trust segregation locks the society into a structurally entrenched state of polarization.

The contributions of this paper are threefold. First, we develop the Dynamic-BABE model, a novel ABM that bridges cognitive psychology, relational network dynamics, and platform economics. Second, we empirically demonstrate the existence of the Polarization Paradox, illustrating the severe trade-off platforms face between short-term financial sustainability and long-term social cohesion. Finally, we highlight the dangers of purely suppressive, content-centric algorithmic moderation, arguing that meaningful solutions to online polarization must move beyond simple avoidance mechanisms toward interventions that foster relational trust and mutual tolerance.
